% ---------------------------------------------------------------------------
% Section: Phase 13 -- parallel and concurrent inference on four CPU cores
% ---------------------------------------------------------------------------
\section{Phase 13: Parallel and Concurrent Inference on Four CPU Cores}
\label{sec:inference}

Phase~13 (30--31 October 2026) asked how far the four cores can be pushed by running work in parallel. The assessment of the
100-prompt run (Section~\ref{sec:impl}) is the starting point: the gate and the exact checks are not the cost (under 2.5\,\% of
model time); decode is, with System~2's thinking at 77.6\,\% and drafts at 17.9\,\%, and six prompts (failed code and timeouts)
used 52\,\%. The research note is \texttt{docs/research/parallel\_inference.md} and the decision is
\texttt{docs/decisions/007-parallel-inference.md}.

\subsection{Literature and what transfers}

Repeated sampling with an exact verifier raises coverage log-linearly with the number of samples
\cite{brown2024monkeys}; hedged requests send a request to several replicas, keep the first answer and cancel the rest
\cite{dean2013tail}; parallel streams over shared parameters behave like a larger model at a fraction of the cost of parameter
scaling \cite{chen2025parscale}; confidence can cut parallel traces early \cite{fu2025deepconf}; instances can attend to a
shared cache \cite{rodionov2025hogwild}; model-free speculation drafts from earlier text \cite{oliaro2024suffix}; table-lookup
kernels speed low-bit models on CPUs \cite{wei2024tmac}. The first two need only what \Sys{} has (an exact verifier and a
batching server) and were implemented. The others need trained transformations, models that cooperate (small ones are reported
not to), or a different runtime, and were not.

\subsection{The measurement that decides: batching and context depth}

llama-server decodes its slots in one batch. Table~\ref{tab:batch} gives the throughput of $N$ concurrent streams on Qwen3-4B-Thinking
(decode only, prompt cached) at three context depths, and Qwen3-1.7B at short context.

\begin{table}[t]
\centering
\caption{Batched decode, $N=4$ streams against one (reference machine, idle).}
\label{tab:batch}
\scriptsize
\begin{tabular}{@{}lrrrr@{}}
\toprule
Model, context & 1 stream & 4 streams & gain & $c(4)$ \\
\midrule
System 1, short & 15.8 & 36.9 & 2.34$\times$ & 1.71 \\
System 2, short & 5.12 & 13.78 & 2.69$\times$ & 1.60 \\
System 2, 800 tokens & 5.03 & 9.17 & 1.82$\times$ & 2.24 \\
System 2, 1,600 tokens & 4.09 & 6.14 & 1.50$\times$ & 3.40 \\
\bottomrule
\end{tabular}
\end{table}

One stream is bandwidth-bound (about 16\,GB/s of weights per second at both sizes), so a few more streams are cheap, until
compute saturates; but the gain falls with context, from 2.7$\times$ to 1.5$\times$ for System~2, and a repair prompt (the earlier
candidate and the verifier's evidence) is about 1,600 tokens deep. A quantised KV cache with flash attention does not restore it
(7.94 against 9.17 tok/s at 800 tokens, 6.36 against 6.14 at 1,600). Self-speculation does not help either: four llama.cpp
n-gram modes gained at most 1.12$\times$ (code repair, 71\,\% of drafted tokens accepted) and lost up to 16\,\% (JSON), because
verifying $k$ drafted tokens costs about $k$ steps on a compute-bound CPU.

\subsection{One request with $n$ choices}

Candidates sent as separate requests make every slot prefill the whole prompt, and prefill (31 tok/s on System~2) is
compute-bound and gains nothing from batching. With a 703-token prompt and a 24-token answer one request takes 31.7\,s, four separate
concurrent requests 124.5\,s ($3.9\times$) and one request with $n=4$ 39.1\,s ($1.23\times$): llama-server prefills once and forks
the KV cache. \texttt{Backend::stream\_n} sends a round as one such request and delivers each choice when its own stream ends, so
it can be verified while the others decode. The first live hedging runs, which used separate requests, were slow for exactly this reason.

\subsection{Fusing the agreement samples with the draft}

For an unverified maths draft the gate always draws $k=2$ cheap samples (A-060), after the draft: $1+c(2)=2.33$ draft times.
Drawn in the draft's batch they cost $c(3)=1.78$, and they are dropped when the draft is stopped early (it is going to System~2
whatever they say). In an interleaved A/B (14 GSM8K prompts, two repetitions, the stream closed at the gate's decision so no
System~2 run is started) the median time to the decision fell from 23.1\,s to 17.9\,s ($-24\,\%$), with the fused arm faster in 23
of 28 pairs (exact sign test $p=0.0005$). Escalations were 7 and 5 of 28 (discordant pairs 4 to 2, McNemar $p=0.69$), and the
decision flips between two repeats of the same arm as often (6 of 28) as between arms, so the difference is sampling noise.

\subsection{When does a concurrent round pay?}

Take $R+1$ rounds, one-attempt success probability $p$ and batch cost factor $c(N)$, and let candidates be independent. A round
of $N$ succeeds with $1-(1-p)^N$, so
\begin{equation}
U(N)=\frac{1-(1-p)^{N(R+1)}}{c(N)\sum_{k=0}^{R}(1-p)^{Nk}}
\label{eq:util}
\end{equation}
is the verified answers per single-attempt time. Widening to $N$ pays iff $U(N)>U(1)$.

\begin{proposition}[Limits of widening] \textnormal{[exact]}
\label{prop:wide}
As $p\to0$, $U(N)/U(1)\to N/c(N)$: widening pays iff the batch is cheaper than the same streams in series. As $p\to1$,
$U(N)/U(1)\to 1/c(N)<1$ for any $c(N)>1$: it never pays.
\end{proposition}
\begin{proof}
For small $p$, $1-(1-p)^{N(R+1)}\approx N(R+1)p$ and $\sum_k(1-p)^{Nk}\approx R+1$, so $U(N)\approx Np/c(N)$ and $U(1)\approx p$.
For $p\to1$ the numerator tends to 1 for every $N$, and the sum to 1, so $U(N)\to1/c(N)$. \qed
\end{proof}

\emph{Consequence.} With the short-context costs the policy widens for $p\le0.48$; with the 1,600-token costs
$[1.0,2.4,3.1,3.4]$ it widens only for $p\le0.07$ (hopeless tasks), because $c(4)=3.4$ is close to 4. The 100-prompt
run also shows that attempts are not independent: of the eight code items that reached System~2, four passed on the first attempt
and four failed all three, and no repair rescued one.

\subsection{Hedging on System 2: implemented, measured, off}

\texttt{nafs\_core::parallel} and the System~2 loop implement it: one future per candidate that generates, applies a patch and
verifies, so verification overlaps the other candidates' decoding; candidates consumed in completion order; the first that passes
wins and dropping the rest cancels their streams; the width from Eq.~(\ref{eq:util}) with an online Beta estimate of $p$; a
\texttt{repair\_rounds} override; \texttt{hedge\_after\_failure}, which runs the first attempt alone; and a round deadline 5\,s inside the
request cap. Table~\ref{tab:hedge} is the live A/B on those eight items (System~2 forced, idle host).

\begin{table}[t]
\centering
\caption{Eight hard code items, System~2 forced (wall time in s; the sequential loop was run twice).}
\label{tab:hedge}
\scriptsize
\begin{tabular}{@{}p{3.3cm}rr@{}}
\toprule
Arm & Verified & Wall \\
\midrule
sequential (1 + 2 serial repairs), run 1 & 4/8 & 4,052 \\
sequential, run 2 & 3/8 & 4,523 \\
four fresh at once, separate requests & 4/8 & 3,308 \\
first alone, then four, one $n=4$ request & 4/7 & 3,558 \\
\bottomrule
\end{tabular}
\end{table}

Four concurrent fresh candidates cost the easy items $+44$ to $+87\,\%$ and saved only on failing ones; the one-request variant was no
faster on the failing items (872--930\,s against 636--901\,s), because a round at repair depth costs 2.2 to 3.4 single attempts
(Table~\ref{tab:batch}). The two sequential runs disagree with each other (4/8 and 3/8), so no accuracy difference can be claimed
at this size. Hedging on System~2 is therefore \emph{not} adopted: \texttt{deliberation.parallel.enabled = false}. Running it also exposed
three defects, all fixed: separate requests prefilling $W$ times; the next attempt's budget sized from the single-stream speed,
which let a round overrun the request cap (HTTP 504); and an error in a round's last candidate discarding the best candidate so far.

\subsection{Cheap concurrent samples before System~2}

Where batching is cheap (System~1, short context) and the verifier is exact, the cores are better spent \emph{before} the
expensive model. When a draft fails its exact check, \Sys{} now draws $k=4$ System~1 samples in one request, verifies each as it finishes and
keeps the first that passes. Let $f$ be the share of requests whose draft fails, $r$ the share of those that a sample rescues,
$\rho$ the cost of the rescue stage (paid on every failed draft) and $E$ the cost of an escalation.

\begin{proposition}[Rescue break-even] \textnormal{[exact]}
\label{prop:rescue}
The saving per request is $f\,(rE-\rho)$, positive iff $r>\rho/E$.
\end{proposition}
\begin{proof}
Of the fraction $f$ with a failed draft, a share $r$ avoids the escalation ($-E$) and all pay $\rho$. \qed
\end{proof}

Offline on 60 MBPP+ items the draft passes 46.7\,\% of the time; of 32 failed drafts, 0, 0, 4 and 6 are rescued by 1, 2, 3 and 4
samples (19\,\% with four, pass@5 56.7\,\%), and four samples in one request take 17.5\,s against 10.4\,s for the draft. Live, in an
interleaved A/B of 40 items (rescue off against on, stream closed at the decision), escalations fell from 22 to 20 (three avoided,
one in the other direction, since the draft itself is random); the rescue added a median 14.5\,s when it ran and 317\,s in all.
Break-even is $317/3=106$\,s per avoided escalation, against 250--580\,s measured for an escalation (the eight code items of the
100-prompt run: median 289\,s, mean 580\,s), so the rescue nets $+10.8$ to $+35.6$\,s per code item and trades no accuracy: every
rescued answer passed the exact check. The sample is small and the two studies agree on the rate; it is on by default
(\texttt{gate.rescue\_samples = 4}), and Lemma~1's call bound gains that term.

\subsection{Serving topology and summary}

System~2 now has four slots and 24,576 tokens of context (6,144 per slot) where it had two, since the 100-prompt run lost 476\,s
to waiting for a System~2 slot. Table~\ref{tab:p13} summarises what was adopted and what was not. Every suite was re-run on the final code (31 October 2026): 206 Rust tests with real Lean, security 26/26, end-to-end 65/65, Rooms 59/59, the harness sweep 50/50 (it now covers the rescue, the parallel settings and their validation), and the nine acceptance tests, where Lemma~1's call bound for the maths request is $1+k+\text{rescue}+N(R+1)=10$ and 8 calls were used.

\begin{table}[t]
\centering
\caption{Phase~13 outcomes.}
\label{tab:p13}
\scriptsize
\begin{tabular}{@{}p{3.0cm}p{1.6cm}p{2.1cm}@{}}
\toprule
Change & Default & Evidence \\
\midrule
Fused agreement samples & on & $-24\,\%$, $p=0.0005$ \\
System~1 rescue (4 samples) & on & 3 of 19 avoided; break-even 106\,s \\
System~2 with 4 slots & on & 2.7$\times$ to 1.5$\times$ decode \\
Concurrent hedging, System~2 & off & no width pays at depth \\
$n$-choice requests & used & 39.1 against 124.5\,s \\
n-gram speculation & no & $\le1.12\times$, down to $0.84\times$ \\
q8\_0 KV cache & no & no gain at depth \\
\bottomrule
\end{tabular}
\end{table}
